\documentclass[11pt]{article}
\usepackage[margin=0.8in]{geometry}
\usepackage{amsmath,microtype}
\newcommand{\guideheading}[1]{\par\medskip\noindent{\large\bfseries #1}\par\smallskip}
\begin{document}
\begin{center}
{\Large Guide: Repairing the Cycle Mismatch with an Inverse Path}\par\medskip
Collatz proof project\quad---\quad September 5, 2026
\end{center}
\guideheading{What the earlier bound ruled out}
The forward bridge search could get stuck after both base orbits reached
the trivial cycle. An integer quantity combining slope, clock, and cycle
phase could only move in one direction. Adding more forward bridges could
not repair the opposite mismatch.

\guideheading{The new route}
We can transport convergence backward along a verified orbit segment.
The number 47 reaches 2 in 65 shortcut steps. Its auxiliary bridge is
$47\mapsto425$, and 425 reaches 2 in 38 steps. The combined route is
\[
 2\ \longleftarrow\ 47\ \longmapsto\ 425\ \longrightarrow\ 2.
\]
Reading the first segment backward makes this a different operation from
the forward paths covered by the earlier bound. Its net change in the
cycle quantity is negative seven. Lean checks the exact orbit computations
and their extension to whole arithmetic progressions.

\guideheading{The premise is still required}
The middle bridge invokes transfer at a smaller parameter. Lean retains
that premise explicitly. The numerical fact that 47 and 425 converge does
not by itself prove convergence for every point on their lifted progressions.
The construction checks that each recursive parameter is smaller for every
quotient, using enough waiting steps when needed.

\guideheading{Why two excursions are enough in principle}
The previous cycle excursion changes the quantity by positive one. The
new one changes it by negative seven. Combining them can correct any
integer mismatch. This gives a construction for a fixed parameter above
15 when both of its base partners are already known to converge.

Moving backward also changes the clock used to define the arithmetic
progression. The binary modulus must retain enough precision for the
largest clock ever visited, including clocks before a backward move.
The program records and checks this condition.

\guideheading{What was tested}
The constructor succeeded for every parameter from 16 through 1000:
985 conditional certificates, with maximum final depth 412. Tests replay
all of them at three quotients, including every reverse segment and every
smaller-parameter comparison. The exact inverse macro is checked in Lean;
the general scheduling argument and the entire census are not yet Lean
coverage theorems.

\guideheading{Why this still does not prove Collatz}
The construction needs the base target to reach the cycle before it can
measure the mismatch. Applying it to every parameter would therefore
require the convergence statement we are trying to prove. The inverse
route fixes a limitation of the certificate language; it does not remove
this global coverage gap.

\guideheading{Verification files}
\begingroup\raggedright
Build \texttt{Collatz.InverseCycleBridge} from \texttt{lean/}.
Run \texttt{python3 analysis/inverse\_cycle\_bridge.py} and
\texttt{python3 -m unittest discover -s analysis -p test\_inverse\_cycle\_bridge.py}.
\par\endgroup
\end{document}
